% Section 1.4, from lectures/L01/appendix/c_exercises.md. The worked solutions are not
% printed; the aside says where they are.

\section{Exercises}
\label{c1:sec:exercises}

Eight, ending with the Cross-check. Do them in order; the last one needs the code from the two before it.

\begin{aside}
\textbf{How to check your work.} A \kindtag{Code} exercise is judged by the suite that ships with it: run \code{make test} at the root of the companion repository, with \code{AEL_DIR} pointing at your toolkit. A suite you have not started is reported as \code{SKIP} rather than as a failure, so the command is worth running from the first day. The hand calculations and designs are checked against the toolkit functions each one names.

\textbf{The solutions are not in this book.} They are published in full in the companion repository, in \file{lectures/L01/appendix}, including the plausible wrong answer where there is one. Read them once you have your own answers, including the ones you are unsure about, because being unsure and right is a different thing from being unsure and wrong and the solutions say which is which.
\end{aside}

% ------------------------------------------------------------------
\recall{The convention}
\label{c1:ex:1}
State the sign convention this book fixes for a current source and for a voltage source's reported current. Then answer this: your solver comes back with a divider output of $-1.71$ V instead of $+1.71$ V, and every voltage in the circuit is negated. Which of the two conventions did you get backwards, and how do you know it was that one rather than the other?

% ------------------------------------------------------------------
\recall{Two dividers}
\label{c1:ex:2}
A divider made of two 1 megohm resistors and a divider made of two 10 ohm resistors have the same ratio.

\begin{enumerate}
  \item What is the Thevenin resistance of each?
  \item One of them is the front end of an oscilloscope probe and the other is inside a power supply. Which is which, and why?
  \item A third divider uses two 10 ohm resistors across a 10 V supply. What is its power dissipation, and is that a problem?
\end{enumerate}

% ------------------------------------------------------------------
\handcalc{The divider, three ways}
\label{c1:ex:3}
For a 33 kilohm over 6.8 kilohm divider on a 10 V supply, by hand:

\begin{enumerate}
  \item The unloaded output.
  \item The Thevenin resistance.
  \item The output with a 10 kilohm load, computed twice: once by substituting the parallel combination into the divider formula, and once from the Thevenin equivalent. Confirm they agree, and say why they must.
  \item The load that would halve the output.
\end{enumerate}
\checkyourself{\code{divider}, \code{divider_output_resistance}}

% ------------------------------------------------------------------
\handcalc{Node equations}
\label{c1:ex:4}
A 1 mA current source injects into node 1. A 1 kilohm resistor runs from node 1 to ground, a 2.2 kilohm resistor from node 1 to node 2, and a 4.7 kilohm resistor from node 2 to ground.

\begin{enumerate}
  \item Write the current-law equation at node 1 and at node 2. Do not solve them yet.
  \item Write the two equations as a matrix, and identify each of the four entries as a stamp contribution from a named resistor.
  \item Solve for both node voltages.
  \item Kill the source and find the resistance looking into node 2. Now divide node 2's voltage from part 3 by the 1 mA source current. Those two numbers are both resistances and both about node 2. Say whether they agree, and if not, say precisely what each one is a resistance \emph{between}.
\end{enumerate}

% ------------------------------------------------------------------
\design{A reference divider}
\label{c1:ex:5}
Design a divider giving 2.5 V from a 10 V supply, subject to:

\begin{itemize}
  \item Thevenin resistance no greater than 2.2 kilohm, so the next stage does not load it.
  \item Divider current no greater than 1 mA, because the supply is a battery.
  \item E12 values only.
\end{itemize}
Give both resistors, the output voltage you actually achieve, the Thevenin resistance and the current. Then state the percentage error in the output, and say whether you would accept it.

There is a tempting wrong answer to this one. Find the pair whose \emph{ratio} is closest to correct, check it against both constraints, and say what it violates.

\checkyourself{\code{divider}, \code{divider_output_resistance}, \code{nearest_e12}}

% ------------------------------------------------------------------
\codeexercise{The netlist}
\label{c1:ex:6}
Implement \code{ael::net::Netlist} to the specification in \secref{c1:sec:b5}.

Four of this chapter's seventeen tests are for this. Getting them green needs no solver, so do this first and confirm the suite goes from four passing to eight.

% ------------------------------------------------------------------
\codeexercise{The solver}
\label{c1:ex:7}
Implement \code{ael::mna::solve} to the same specification.

Assemble the stamps, solve by Gaussian elimination with partial pivoting, and detect a singular matrix rather than returning zeros. All seventeen tests should pass.

Two of them are worth reading before you start rather than after you fail them: \code{Mna.ParallelResistorsCombine}, which fails for a solver that indexes by element instead of by node, and \code{Mna.FloatingNodeIsReportedRatherThanGuessed}, which fails for a solver that trusts its own elimination.

% ------------------------------------------------------------------
\crosscheck{The loaded divider}
\label{c1:ex:8}
The signature exercise of this book, and the shape every later one takes. Compute the same number three ways and reconcile them.

Take the 33 kilohm over 6.8 kilohm divider on a 10 V supply, loaded with 10 kilohm.

\begin{enumerate}
  \item \textbf{By hand.} Substitute the parallel combination into the divider formula. Write down the answer to three significant figures before doing anything else.
  \item \textbf{By your closed form.} Whatever you wrote for Exercise~\ref{c1:ex:3}, as a function, evaluated on the same numbers.
  \item \textbf{By your solver.} Build the netlist with four elements, one voltage source and three resistors, and read node 1 out of \code{ael::mna::solve}.
\end{enumerate}
Then reconcile. State the difference between each pair, and say what would have to be true for each difference to be non-zero.

\task{What to expect}
\textbf{Legs 1 and 2 should agree to every digit you wrote down.} They are the same formula evaluated by two different pieces of hardware. A disagreement here is an arithmetic slip in leg 1 or a typo in leg 2, and it is not interesting except that finding it is faster than finding it later.

\textbf{Legs 2 and 3 should agree to about twelve significant figures.} Both are solving the same linear problem in double precision, one by algebra you did in advance and one by elimination. The difference is rounding and nothing else, and it should be somewhere around $10^{-13}$ relative.

\textbf{A difference of a few per cent means one of them is not solving the circuit you think it is.} The usual cause is the load stamped between the wrong pair of nodes, which quietly gives you a different circuit rather than an error. The second usual cause is the supply node and the output node swapped, which gives you the reciprocal divider.

\textbf{A difference of exactly a factor of two, or exactly the unloaded answer, means the load is not in the circuit at all.} Check \code{resistorCount()}.

This exercise is deliberately one where all three legs agree. That is worth doing once, so that you know what agreement looks like before Chapter~\ref{c7:ch:smallsignal}, where a closed form and the solver disagree by a factor of ten and the closed form is the one that is wrong.
